
\subsection{State Management and Persistence}

All workflow data is stored in a single \texttt{MasterState} object defined in \texttt{state.py}. This Pydantic-based state class contains over 50 typed attributes covering:

\begin{itemize}
    \item \textbf{User Input}: Original message, intent classification, routing metadata
    \item \textbf{Firmware Workflow}: Board name, MCU series, .ioc file path, CubeMX output directory
    \item \textbf{AI Workflow}: Model path, task type, target MCU, compression level, search iteration count
    \item \textbf{Customization}: Model layers, modification parameters, NNI experiment directory
    \item \textbf{Integration}: Generated code paths, build status
    \item \textbf{Web Research}: Query history, search results, evaluation scores
\end{itemize}

This persistent state model enables \textbf{Context Chaining}, a key architectural feature where outputs from upstream workflows (e.g., the specific MCU target selected during \texttt{firmware\_flow}) automatically populate the input context for downstream workflows (\texttt{ai\_flow}). This eliminates redundant user data entry and prevents configuration mismatches, such as selecting an AI model incompatible with the previously generated firmware board. The state object acts as a \textit{single source of truth}, synchronizing project parameters across heterogeneous toolchains (CubeMX, STEdgeAI, TensorFlow).

LangGraph automatically persists this state at every node transition through a checkpointing mechanism, enabling:

\begin{enumerate}
    \item \textbf{Fault Tolerance}: Resume execution from the last successful checkpoint after transient failures
    \item \textbf{Human-in-the-Loop Gates}: Suspend execution for user approval without losing context
    \item \textbf{Time-Travel Debugging}: Replay past executions to diagnose workflow errors
\end{enumerate}

A detailed analysis of the underlying storage infrastructure and the dual-layer memory management logic is provided in Section \ref{sec:redis_architecture}.

\subsection{Reliability in LLM-Driven Orchestration}

Relying on Large Language Models for orchestration decisions introduces non-determinism. During stress testing with concurrent users, unconstrained LLM outputs occasionally caused validation crashes or incorrect routing. The framework solves this by enforcing three layers of input normalization before any LLM response modifies the \texttt{MasterState}.

First, I implemented type coercion for numerical parameters. When extracting data like sample rates for synthetic data generation, LLMs frequently append human-readable units (e.g., \textit{``2 kHz''} instead of \texttt{2000}). The framework preemptively normalizes these strings into raw floating-point measurements before passing them to Pydantic, bypassing strict validation errors.

Second, I introduced regex-based fallbacks for hardware identification. If the LLM extracts an imprecisely formatted identifier (hallucinating prefixes like \textit{``STM32 F401''} instead of \texttt{STM32F401}), a relaxed regular expression such as:

\texttt{r'(?:STM32)?[\textbackslash{}s-]*([A-Z])([0-9])'}

\noindent guarantees the core MCU series letter and digit are captured correctly for the downstream ST toolchains.

Finally, I adopted an enumerated prompt anchoring strategy. During model discovery, providing the LLM with only the total count of available models resulted in hallucinated index selections when users specified model names. Injecting the complete, enumerated list of model names directly into the extraction prompt creates a bounded decision space and enforces deterministic routing.

\subsection{System Integrity and Transactional Registration}

To maintain the stability of the model registry, the framework implements a \textbf{Transactional Registration Flow} for user-provided models. When a user submits a new model URL, the entry is first held in a \textit{Pending} state within the \texttt{MasterState}. The system then performs an autonomous validation via the STEdgeAI analysis engine. Only upon successful technical profiling (\texttt{analyze\_success = true}) is the model permanently committed to the \texttt{predefined\_models.json} registry. This transactional behavior ensures that the system's knowledge base remains free from broken links, malformed architectures, or incompatible file formats.

\subsection{Configuration Management}

Runtime parameters are managed through a dedicated \texttt{Configuration} class (\texttt{configuration.\allowbreak py}) that dynamically loads settings from a prioritized hierarchy of sources. \textbf{Environment Variables} are used for sensitive data such as STMicroelectronics credentials (\texttt{ST\_EMAIL}, \texttt{ST\_PASSWORD}) and API keys, while a local \textbf{Config File} (\texttt{config.json}) defines system-specific paths including the STM32CubeMX binary location, the STM32Cube repository path, and AI output directories. 

In addition, the \textbf{LangGraph Runtime Config} enables the injection of per-request settings directly into each workflow node via the \texttt{config} parameter.

The framework relies on several important parameters to govern its behavior: \texttt{local\_llm} defines the language model identifier (defaulting to \texttt{mistral:latest}), while \texttt{llm\_temperature} balances deterministic routing ($0.0$) with creative code generation ($0.2$). The \texttt{llm\_\allowbreak context\_\allowbreak window} manages token limits, typically $4096$--$8192$, and \texttt{cubemx\_path} points to the absolute location of the STM32CubeMX CLI. 

For AI-specific workflows, the \texttt{ai\_target} specifies the target MCU (e.g., \texttt{stm32h743}), and \texttt{ai\_compression} sets the desired quantization level ranging from \texttt{low} to \texttt{very\_high}. All parameters are validated at runtime using \textbf{Pydantic validators}, to ensure that required paths exist and credentials are correctly formatted before the orchestrator starts any workflow execution.

\subsection{Graph Topology and Node Composition}

Each of the five primary workflows is implemented as a \textbf{subgraph} with its own internal routing logic. The master graph connects these subgraphs through conditional edges that allow \textit{inter-workflow transitions} based on user decisions:

\begin{itemize}
    \item After completing firmware generation, the system offers to proceed to AI model discovery
    \item After AI model optimization, the system offers to integrate the generated C code into the firmware project
    \item At any point, users can invoke web research to gather technical documentation
\end{itemize}

This multi-stage pipeline mirrors the typical embedded AI development lifecycle:

\begin{center}
    \texttt{Hardware Setup → Model Selection → Model Customization → Code Generation → Integration → Testing}
\end{center}

\noindent Figure~\ref{fig:workflow_graph} illustrates the full graph topology, including the feedback cycles that distinguish this architecture from a simple linear pipeline.

% Five Primary Workflows — Interconnected Graph Diagram
\begin{figure}[htbp]
    \centering
    \begin{tikzpicture}[
            node distance=1.4cm and 2.2cm,
            every node/.style={font=\sffamily\small},
            wf/.style={draw, rounded corners, minimum width=3.4cm, minimum height=1cm, align=center, text=black, inner sep=0.15cm},
            decision/.style={draw, diamond, fill=gray!10, minimum width=1.2cm, minimum height=0.8cm, align=center, inner sep=1pt, aspect=2.2, font=\sffamily\scriptsize},
            arrow/.style={->, thick, >=stealth},
            cycle/.style={->, thick, >=stealth, dashed, bend left=25},
            label/.style={font=\sffamily\scriptsize, text=black}
        ]

        % === ROW 1 (left to right): Router → Firmware ===
        \node[wf, fill=gray!15] (router) {\textbf{LLM Router}\\(Request Classification)};
        \node[wf, fill=green!15, right=2.5cm of router] (firmware) {\textbf{1. Firmware Generation}\\(\texttt{firmware\_flow})};

        % === ROW 2 (right to left): AI Discovery ← Decision ===
        \node[decision, below=1.2cm of firmware] (d1) {Continue\\to AI?};
        \node[wf, fill=blue!12, left=2.5cm of d1] (ai) {\textbf{2. AI Model Discovery}\\(\texttt{ai\_flow})};

        % === ROW 3 (left to right): Customization ===
        \node[wf, fill=purple!12, below=1.2cm of ai] (custom) {\textbf{3. Customization}\\(\texttt{customization\_flow})};
        \node[decision, right=2.5cm of custom] (d2) {Continue to\\Integration?};

        % === ROW 4 (right to left): Integration ===
        \node[wf, fill=cyan!15, below=1.2cm of d2] (integration) {\textbf{4. Integration}\\(\texttt{integration\_flow})};

        % === Web Search (independent branch, left side) ===
        \node[wf, fill=orange!15, below=1.2cm of router] (search) {\textbf{5. Web Research}\\(\texttt{search\_flow})};

        % === END node ===
        \node[draw, rounded corners, fill=red!10, minimum width=1.5cm, minimum height=0.6cm, left=2.5cm of integration] (endnode) {\textbf{END}};

        % ============================================================
        % MAIN FLOW ARROWS
        % ============================================================

        % Router dispatches
        \draw[arrow] (router) -- node[label, above] {route: firmware} (firmware);
        \draw[arrow] (router) -- node[label, left] {route: search} (search);
        \draw[arrow] (router.south) -- ++(0,-0.4) -| node[label, near start, right] {route: ai} (ai.north);

        % Firmware → Decision 1
        \draw[arrow] (firmware) -- (d1);

        % Decision 1 branches
        \draw[arrow] (d1) -- node[label, above] {Yes} (ai);
        \draw[arrow] (d1.east) -- ++(0.8,0) node[label, right] {No → END};

        % AI → Customization
        \draw[arrow] (ai) -- (custom);

        % Customization → Decision 2
        \draw[arrow] (custom) -- (d2);

        % Decision 2 branches
        \draw[arrow] (d2) -- node[label, right] {Yes} (integration);
        \draw[arrow] (d2.east) -- ++(0.8,0) node[label, right] {No → END};

        % Integration → END
        \draw[arrow] (integration) -- (endnode);

        % Search → END
        \draw[arrow] (search.south) -- ++(0,-0.4) -| (endnode.north);

        % ============================================================
        % CYCLES (graph nature, not tree)
        % ============================================================

        % Resource failure: AI Analysis → back to AI Discovery (retry model)
        \draw[cycle, color=red!60] (custom.west) -- ++(-1.2,0) |- node[label, left, near start, text=red!60] {\scriptsize model too large} (ai.west);

        % Resource failure: change board → back to Firmware
        \draw[cycle, color=red!60] (ai.north) -- ++(0, 0.5) -| node[label, above, near start, text=red!60] {\scriptsize change board} (firmware.north);

        % Router can also dispatch directly to Integration
        \draw[arrow, dotted, color=gray!70] (router.south) -- ++(0,-0.4) -| (integration.west);

        % Human decision after customization can loop back for more modifications
        \draw[cycle, color=blue!50] (custom.south) -- ++(0,-0.4) -- ++(-2.2,0) -- ++(0,1.8) -- node[label, left, text=blue!50] {\scriptsize re-modify} (custom.north west);

    \end{tikzpicture}
    \caption{Interconnected graph of the five primary workflows. Solid arrows represent the main execution pipeline, dashed arrows denote feedback cycles triggered by resource constraint violations or user-driven re-modifications. Dotted arrows indicate direct routing alternatives (bypassing upstream workflows).}
    \label{fig:workflow_graph}
\end{figure}

The graph construction logic resides in four primary ``injection'' functions:

\begin{itemize}
    \item \texttt{inject\_\allowbreak firmware\_\allowbreak nodes()}: \textit{Firmware Generation} workflow
    \item \texttt{inject\_\allowbreak ai\_\allowbreak analysis\_\allowbreak nodes()}: \textit{AI Model Discovery}, \textit{Model Customization}, and data management workflows
    \item \texttt{inject\_\allowbreak integration\_\allowbreak nodes()}: \textit{Firmware-AI Integration} workflow
    \item \texttt{inject\_\allowbreak web\_\allowbreak search\_\allowbreak nodes()}: \textit{Web Research} workflow
\end{itemize}

Each function encapsulates the node definitions and internal routing logic for its respective domain, adding them to the master \texttt{StateGraph} builder.

\subsection{Human-in-the-Loop Interaction}
\label{sec:hitl_interaction}

The framework places humans at critical decision points using LangGraph's \texttt{interrupt()} primitive. This mechanism pauses execution and returns control to the API server until the user provides feedback.

\subsubsection{Interrupt and Resume Mechanism}

LangGraph snapshots the full graph state after every node transition. When a node calls \texttt{interrupt()}, the runtime suspends the flow and surfaces a payload to the client. The code below illustrates an approval gate between the firmware and AI workflows:

\begin{lstlisting}[language=Python, caption={Interrupt-based approval gate}]
# From workflow1_firmware.py
def decide_continue_to_ai(state, config):
    # Suspend execution and request user feedback
    user_response = interrupt({
        "type": "decision_request",
        "message": "Firmware generated. Continue with AI selection?",
        "options": ["yes", "no", "skip"]
    })
    
    # Analyze response via LLM
    llm_parser = get_llm(config, structured_schema=DecisionExtraction, temperature=0)
    result = llm_parser.invoke([
        SystemMessage(content=instructions),
        HumanMessage(content=user_response)
    ])
    
    state.route = "ai_flow" if result.continue_to_ai else "END"
    return state
\end{lstlisting}

When the user replies, the graph resumes from the exact checkpoint. Since Redis handles the underlying state storage (as detailed in Section \ref{sec:redis_architecture}), even multi-hour pauses do not affect system stability.

\subsubsection{Interaction Patterns}

We categorize human interventions into two patterns. **Approval Gates** handle high-risk transitions, such as confirming hardware parameters or accepting a candidate model during iterative search. **Natural Language Steering** instead allows the user to redirect the agent mid-stream. For example, a user might state, \textit{``I actually need a smaller model''}, triggering a transition back to the search node with updated constraints.

\subsubsection{Actionable Reporting}

To prevent information overload, the orchestrator parses raw logs from tools like CubeMX or STEdgeAI. Instead of displaying hundreds of lines, the system extracts a one-line summary, such as \textit{``Project generated: STM32H7, 4 peripherals enabled''} or \textit{``Memory Warning: Flash limit exceeded''}. These summaries populate the context-aware prompts at interrupt points, ensuring the user makes informed decisions without reviewing verbose technical output.

\subsection{Logging and Observability}

The system implements a multi-tier observability strategy, combining structured backend logging with real-time progress streaming. Technical logs are managed via the standard Python \texttt{logging} library with distinct severity levels.

\textbf{DEBUG} logs provide detailed execution traces, including LLM prompt-response pairs and internal state transitions, while \textbf{INFO} logs track high-level operational events such as workflow entry, node completion, and subprocess execution. \textbf{WARNING} logs are reserved for non-fatal anomalies, such as low LLM confidence scores or active fallback mechanisms, and \textbf{ERROR} logs capture recoverable failures, including API timeouts or malformed user inputs requiring re-interaction.

A specific architectural feature is the log interception mechanism: logs generated by critical nodes, such as training status or hardware configuration, are captured by a dedicated \texttt{QueueHandler} and streamed to the client using the \textbf{Newline Delimited JSON (NDJSON)} format. This allows frontend clients, such as the \textit{Gradio UI} or the \textbf{VS Code extension}, to provide immediate status visualization and progress bars, transforming a complex multi-step backend process into a transparent and interactive user experience.
